\documentclass[11pt]{article}
\usepackage{fontspec}
\setmainfont{Latin Modern Roman}
\usepackage{amsmath,amssymb,mathtools,unicode-math}
\setmathfont{Latin Modern Math}
\usepackage{xurl,hyperref}
\hypersetup{hidelinks,pdftitle={Problem statement}}
\newcommand{\sref}[2]{\hyperlink{source:#1:#2}{[S#2]}}
\newcommand{\cataloguescope}[1]{\paragraph{Mathematical question.}#1}
\begin{document}
\section{Decomposing a connected graph into paths.}\label{statement}

\subsection{Definitions and mathematical statement}

\paragraph{Shared notation.}$\mathbb N=\{1,2,\ldots\}$, and $\mathbb Z,\mathbb Q,\mathbb R,\mathbb C$ denote the integers, rationals, reals and complex numbers. Any other conventions are specified locally.


For a finite simple graph $G$, a path is a sequence of distinct vertices with edges between consecutive vertices. A path decomposition is a partition of $E(G)$ into edge sets of paths; vertices may occur in different paths. Write $n=|V(G)|$.
\paragraph{Statement recorded in the supplied source.}
Does every connected graph $G$ admit a path decomposition with at most \[\left\lfloor\frac{n+1}{2}\right\rfloor=\left\lceil\frac n2\right\rceil\] paths? \sref{statement}{2}
\subsection{Short English statement}
Can the edges of every connected graph be partitioned into at most half its order rounded up paths?
\subsection{Sources}
\begin{description}
\item[\hypertarget{source:statement:1}{S1}] ULAM.AI and the UnsolvedMath Contributors, UnsolvedMath: A Curated Collection of Open Mathematics Problems, catalogue entry. CC BY 4.0. \url{https://huggingface.co/datasets/ulamai/UnsolvedMath}.
\item[\hypertarget{source:statement:2}{S2}] Open Problem Garden, Decomposing a connected graph into paths., original node 46583, statement and mathematical discussion. \url{https://www.openproblemgarden.org/op/decomposing_a_connected_graph_into_paths}.
\end{description}
\label{end:statement}
\end{document}
